% Junior Level 2, Week 1 -- answer key
\answersweek{1}{Forever and Bounce}{Answers for workbook chapter 1 and play-at-home card 1.}

\section{Lesson questions}
\begin{answer}[Show-me warm-up]
The Week 1 show-me warm-up is described, with what to look for, in the Assessments booklet. It is ungraded: use it to decide how much help each child needs in the first weeks.
\end{answer}

\begin{answer}[Exit question]
Forever. A repeat 4 does its blocks 4 times and then stops by itself. A forever loop keeps going until someone clicks the red stop sign.
\end{answer}

\begin{answer}[Reflection]
Things that go on until someone stops them: a merry-go-round, a fan, a song on repeat, a video that plays the next one, a washing machine (which actually stops by itself: a good chance to talk about the difference). Either loop the child circles is right; most will circle forever this week.
\end{answer}

\section{Workbook exercises}
\begin{answer}[Exercise 1.1: Bounce Cat]
\textbf{What you should see:} after the green flag, the cat runs to the right edge, turns around, runs to the left edge, turns around again, and keeps going until the red stop sign is clicked. It always faces the way it walks and never goes upside down.

\textbf{What to look for:} \tickbox{} \enquote{point in direction 90} and \enquote{set rotation style left-right} are above the forever \quad \tickbox{} \enquote{move 10 steps} and \enquote{if on edge, bounce} are inside the forever \quad \tickbox{} the child stops it with the stop sign.
\end{answer}

\begin{answer}[Exercise 1.2: Dance Forever]
\textbf{What you should see:} pressing space makes the Ballerina change costume 4 times, one each second, and then stop. Clicking the green flag makes her change costume every second without stopping. The red stop sign stops her.

\textbf{Which one stopped by itself?} The repeat (the space-key script). The forever loop needed the stop sign.

\textbf{What to look for:} \tickbox{} both scripts are in the Ballerina's code, not the cat's \quad \tickbox{} the child can say \enquote{repeat stops, forever keeps going}.
\end{answer}

\begin{answer}[Exercise 1.3: The Missing Hello]
\textbf{Circle:} the say block and the sound block under the forever. They are not snapped on (forever has a flat bottom), and even if they were, a forever loop never finishes, so their turn would never come.

\textbf{Fix:} drag them up and snap them between \enquote{when green flag clicked} and the forever.

\textbf{What you should see after the fix:} the cat says \enquote{Watch me!} for 2 seconds, then meows (the meow starts as the cat begins to move), then bounces from side to side until the stop sign is clicked. If the sound block is snapped above the say block, the cat meows while it talks: also fine.
\end{answer}

\section{Play-at-home card}
\begin{answer}[Forever freeze dance]
The child keeps dancing until the music pauses, and freezes then, not before. In the bounce walk, they turn around when they touch the wall. Ask: \enquote{Who stopped the dance?} (The one with the pause button.)
\end{answer}

\begin{answer}[My stop plan]
Any clear plan is right: \enquote{When the timer rings, I press stop, save, and go outside}, \enquote{I put the tablet on the shelf and get a snack.} Look for a drawing of a timer or a stop sign and a next activity away from the screen.
\end{answer}
